\documentclass[10pt,a4paper]{amsart}
\usepackage{amsmath,amssymb,amsthm,graphicx}
\usepackage{verbatim}
\newenvironment{reference}{\comment}{\endcomment}
%\newenvironment{reference}{}{}
\newenvironment{slogan}{\comment}{\endcomment}
\newenvironment{history}{\comment}{\endcomment}

% For commutative diagrams we use Xy-pic
\usepackage[all]{xy}

% We use 2cell for 2-commutative diagrams.
\xyoption{2cell}
% Theorem environments.
%
\theoremstyle{plain}
\newtheorem{theorem}[subsection]{Theorem}
\newtheorem{proposition}[subsection]{Proposition}
\newtheorem{lemma}[subsection]{Lemma}

\theoremstyle{definition}
\newtheorem{definition}[subsection]{Definition}
\newtheorem{example}[subsection]{Example}
\newtheorem{exercise}[subsection]{Exercise}
\newtheorem{situation}[subsection]{Situation}

\theoremstyle{remark}
\newtheorem{remark}[subsection]{Remark}
\newtheorem{remarks}[subsection]{Remarks}

\numberwithin{equation}{subsection}

% Macros
%
\def\lim{\mathop{\mathrm{lim}}\nolimits}
\def\colim{\mathop{\mathrm{colim}}\nolimits}
\def\Spec{\mathop{\mathrm{Spec}}}
\def\Hom{\mathop{\mathrm{Hom}}\nolimits}
\def\Ext{\mathop{\mathrm{Ext}}\nolimits}
\def\SheafHom{\mathop{\mathcal{H}\!\mathit{om}}\nolimits}
\def\SheafExt{\mathop{\mathcal{E}\!\mathit{xt}}\nolimits}
\def\Sch{\mathit{Sch}}
\def\Mor{\mathop{\mathrm{Mor}}\nolimits}
\def\Ob{\mathop{\mathrm{Ob}}\nolimits}
\def\Sh{\mathop{\mathit{Sh}}\nolimits}
\def\NL{\mathop{N\!L}\nolimits}
\def\CH{\mathop{\mathrm{CH}}\nolimits}
\def\proetale{{pro\text{-}\acute{e}tale}}
\def\etale{{\acute{e}tale}}
\def\QCoh{\mathit{QCoh}}
\def\Ker{\mathop{\mathrm{Ker}}}
\def\Im{\mathop{\mathrm{Im}}}
\def\Coker{\mathop{\mathrm{Coker}}}
\def\Coim{\mathop{\mathrm{Coim}}}

% Boxtimes
%
\DeclareMathSymbol{\boxtimes}{\mathbin}{AMSa}{"02}

%
% Macros for moduli stacks/spaces
%
\def\QCohstack{\mathcal{QC}\!\mathit{oh}}
\def\Cohstack{\mathcal{C}\!\mathit{oh}}
\def\Spacesstack{\mathcal{S}\!\mathit{paces}}
\def\Quotfunctor{\mathrm{Quot}}
\def\Hilbfunctor{\mathrm{Hilb}}
\def\Curvesstack{\mathcal{C}\!\mathit{urves}}
\def\Polarizedstack{\mathcal{P}\!\mathit{olarized}}
\def\Complexesstack{\mathcal{C}\!\mathit{omplexes}}
% \Pic is the operator that assigns to X its picard group, usage \Pic(X)
% \Picardstack_{X/B} denotes the Picard stack of X over B
% \Picardfunctor_{X/B} denotes the Picard functor of X over B
\def\Pic{\mathop{\mathrm{Pic}}\nolimits}
\def\Picardstack{\mathcal{P}\!\mathit{ic}}
\def\Picardfunctor{\mathrm{Pic}}
\def\Deformationcategory{\mathcal{D}\!\mathit{ef}}
\usepackage{cleveref}
\crefname{originallemma}{original lemma}{original lemmas}
\Crefname{originallemma}{Original lemma}{Original lemmas}
\makeatletter
\newlabel{cohomology-lemma-internal-hom-evaluate}{{0A8U}{source}{}{sourcepointer-1}{}}
\newlabel{cohomology-lemma-internal-hom-evaluate@cref}{{[originallemma][0][]0A8U}{[1][0][]source}{}{}{}}
\makeatother
\crefname{originallemma}{original lemma}{original lemmas}
\Crefname{originallemma}{Original lemma}{Original lemmas}
\makeatletter
\newlabel{perfect-lemma-affine-compare-bounded}{{06Z0}{source}{}{sourcepointer-2}{}}
\newlabel{perfect-lemma-affine-compare-bounded@cref}{{[originallemma][0][]06Z0}{[1][0][]source}{}{}{}}
\makeatother
\crefname{originallemma}{original lemma}{original lemmas}
\Crefname{originallemma}{Original lemma}{Original lemmas}
\makeatletter
\newlabel{perfect-lemma-coherent-internal-hom}{{0D0C}{source}{}{sourcepointer-3}{}}
\newlabel{perfect-lemma-coherent-internal-hom@cref}{{[originallemma][0][]0D0C}{[1][0][]source}{}{}{}}
\makeatother
\crefname{originallemma}{original lemma}{original lemmas}
\Crefname{originallemma}{Original lemma}{Original lemmas}
\makeatletter
\newlabel{lemma-affine-duality}{{0A88}{source}{}{sourcepointer-4}{}}
\newlabel{lemma-affine-duality@cref}{{[originallemma][0][]0A88}{[1][0][]source}{}{}{}}
\makeatother
\crefname{originallemma}{original lemma}{original lemmas}
\Crefname{originallemma}{Original lemma}{Original lemmas}
\makeatletter
\newlabel{perfect-lemma-quasi-coherence-internal-hom}{{0A6H}{source}{}{sourcepointer-5}{}}
\newlabel{perfect-lemma-quasi-coherence-internal-hom@cref}{{[originallemma][0][]0A6H}{[1][0][]source}{}{}{}}
\makeatother
\crefname{originallemma}{original lemma}{original lemmas}
\Crefname{originallemma}{Original lemma}{Original lemmas}
\makeatletter
\newlabel{more-algebra-lemma-internal-hom-evaluate-isomorphism}{{0A68}{source}{}{sourcepointer-6}{}}
\newlabel{more-algebra-lemma-internal-hom-evaluate-isomorphism@cref}{{[originallemma][0][]0A68}{[1][0][]source}{}{}{}}
\makeatother
\crefname{originallemma}{original lemma}{original lemmas}
\Crefname{originallemma}{Original lemma}{Original lemmas}
\makeatletter
\newlabel{dualizing-lemma-finite-ext-into-bounded-injective}{{0G4H}{source}{}{sourcepointer-7}{}}
\newlabel{dualizing-lemma-finite-ext-into-bounded-injective@cref}{{[originallemma][0][]0G4H}{[1][0][]source}{}{}{}}
\makeatother
\crefname{originallemma}{original lemma}{original lemmas}
\Crefname{originallemma}{Original lemma}{Original lemmas}
\makeatletter
\newlabel{more-algebra-lemma-internal-hom-evaluate-isomorphism-technical}{{0A69}{source}{}{sourcepointer-8}{}}
\newlabel{more-algebra-lemma-internal-hom-evaluate-isomorphism-technical@cref}{{[originallemma][0][]0A69}{[1][0][]source}{}{}{}}
\makeatother
\title{The internal Hom evaluation map is an isomorphism under boundedness and finite injective dimension hypotheses}
\author{The Stacks Project Authors}
\date{stacks-project:duality.tex:local-snapshot-20261012; GNU FDL 1.2 or later}
\begin{document}
\maketitle
\noindent\textit{Editorial note.} All delivered mathematical text is the author's own, copied from the pinned Stacks Project chapter file duality.tex; no statement or proof step was rewritten, reordered or removed.\par\medskip
\noindent\textit{Editorial note.} Preamble provenance: the pinned chapter file opens with an input of the shared preamble and carries no preamble text of its own, so the project's single Stacks-wide preamble staged at sources/stacks-duality/preamble.tex is excerpted verbatim under the recipe fixed by LANE\_BRIEF 22.1, as three spans, lines 8-13 (verbatim package and the reference/slogan/history environments), lines 14-18 (Xy-pic setup) and lines 154-217 (the native theorem declarations and mathematical macros). Nothing was added, removed or redefined; the wrapper class and the amsmath/amssymb/amsthm/graphicx base packages are supplied by the assembler.\par\medskip
\noindent\textit{Editorial note.} The counted statement and proof contain no cross reference, no citation and no included file, so the delivered body is self contained.\par\medskip
\noindent\textit{Editorial note.} Canonical location of the counted proposition: Stacks tag 0G4I (sources/stacks-duality/pinned-original-source.tex line 415).\par\medskip
\noindent\textit{Allowed standard prerequisites (editorial).}\begin{itemize}
\item The evaluation map between internal Homs. Cohomology, Lemma labelled cohomology-lemma-internal-hom-evaluate, invoked in the statement. Supplies the arrow whose bijectivity is asserted.
\item Comparison of the affine and the scheme level descriptions. Same chapter, Lemma labelled lemma-affine-duality, cited at proof lines 454-458 and 464-466, and Derived Categories of Schemes, Lemma labelled perfect-lemma-quasi-coherence-internal-hom, cited at proof lines 456-460. Shows both sides commute with the functor from D(A) to D\_QCoh(O\_X).
\item Boundedness of the internal Hom of coherent complexes. Derived Categories of Schemes, Lemma labelled perfect-lemma-coherent-internal-hom, cited at proof lines 450-454, and Dualizing Complexes, Lemma labelled dualizing-lemma-finite-ext-into-bounded-injective, cited at proof lines 466-468. Gives RSheafHom(K, L) in D\textasciicircum{}+\_Coh in case (1) and in D\_Coh in case (2).
\item The algebraic evaluation isomorphisms. More on Algebra, Lemmas labelled more-algebra-lemma-internal-hom-evaluate-isomorphism parts (1) and (2), cited at proof lines 460-462 and 468-469. Supply the isomorphism after the reduction to affines.
\end{itemize}\medskip
\ifdefined\hypertarget\hypertarget{sourcepointer-1}{}\fi
\noindent\textit{Source reference.} Original lemma 0A8U: sources/stacks-chow/tags.txt, source lines 12927--12927. Cross-chapter standard prerequisite located in the project tag list staged in the local snapshot; invoked at proof line None of the counted proof. It is retained only as a navigation pointer; every load-bearing step is inside the delivered span.\par\medskip
\ifdefined\hypertarget\hypertarget{sourcepointer-2}{}\fi
\noindent\textit{Source reference.} Original lemma 06Z0: sources/stacks-chow/tags.txt, source lines 8904--8904. Cross-chapter standard prerequisite located in the project tag list staged in the local snapshot; invoked at proof line 440 of the counted proof. It is retained only as a navigation pointer; every load-bearing step is inside the delivered span.\par\medskip
\ifdefined\hypertarget\hypertarget{sourcepointer-3}{}\fi
\noindent\textit{Source reference.} Original lemma 0D0C: sources/stacks-chow/tags.txt, source lines 16307--16307. Cross-chapter standard prerequisite located in the project tag list staged in the local snapshot; invoked at proof line 448 of the counted proof. It is retained only as a navigation pointer; every load-bearing step is inside the delivered span.\par\medskip
\ifdefined\hypertarget\hypertarget{sourcepointer-4}{}\fi
\noindent\textit{Source reference.} Original lemma 0A88: sources/stacks-duality/pinned-original-source.tex, source lines 366--366. Same-chapter pointer into the author environment lemma at source line 366; invoked at proof line 451 of the counted proof. It is retained only as a navigation pointer; every load-bearing step is inside the delivered span.\par\medskip
\ifdefined\hypertarget\hypertarget{sourcepointer-5}{}\fi
\noindent\textit{Source reference.} Original lemma 0A6H: sources/stacks-chow/tags.txt, source lines 12845--12845. Cross-chapter standard prerequisite located in the project tag list staged in the local snapshot; invoked at proof line 453 of the counted proof. It is retained only as a navigation pointer; every load-bearing step is inside the delivered span.\par\medskip
\ifdefined\hypertarget\hypertarget{sourcepointer-6}{}\fi
\noindent\textit{Source reference.} Original lemma 0A68: sources/stacks-chow/tags.txt, source lines 12836--12836. Cross-chapter standard prerequisite located in the project tag list staged in the local snapshot; invoked at proof line 458 of the counted proof. It is retained only as a navigation pointer; every load-bearing step is inside the delivered span.\par\medskip
\ifdefined\hypertarget\hypertarget{sourcepointer-7}{}\fi
\noindent\textit{Source reference.} Original lemma 0G4H: sources/stacks-chow/tags.txt, source lines 20127--20127. Cross-chapter standard prerequisite located in the project tag list staged in the local snapshot; invoked at proof line 465 of the counted proof. It is retained only as a navigation pointer; every load-bearing step is inside the delivered span.\par\medskip
\ifdefined\hypertarget\hypertarget{sourcepointer-8}{}\fi
\noindent\textit{Source reference.} Original lemma 0A69: sources/stacks-chow/tags.txt, source lines 12837--12837. Cross-chapter standard prerequisite located in the project tag list staged in the local snapshot; invoked at proof line 468 of the counted proof. It is retained only as a navigation pointer; every load-bearing step is inside the delivered span.\par\medskip
\begin{lemma}
\label{lemma-internal-hom-evaluate-isom}
Let $X$ be a Noetherian scheme. Let $K, L, M \in D_\QCoh(\mathcal{O}_X)$. 
Then the map
$$
R\SheafHom(L, M) \otimes_{\mathcal{O}_X}^\mathbf{L} K
\longrightarrow
R\SheafHom(R\SheafHom(K, L), M)
$$
of Cohomology, Lemma \ref{cohomology-lemma-internal-hom-evaluate}
is an isomorphism in the following two cases
\begin{enumerate}
\item $K \in D^-_{\textit{Coh}}(\mathcal{O}_X)$,
$L \in D^+_{\textit{Coh}}(\mathcal{O}_X)$, and $M$ affine locally has
finite injective dimension (see proof), or
\item $K$ and $L$ are in $D_{\textit{Coh}}(\mathcal{O}_X)$,
the object $R\SheafHom(L, M)$ has finite tor dimension, and
$L$ and $M$ affine locally have finite injective dimension
(in particular $L$ and $M$ are bounded).
\end{enumerate}
\end{lemma}

\begin{proof}
Proof of (1). We say $M$ has affine locally finite injective dimension
if $X$ has an open covering by affines $U = \Spec(A)$ such that the object
of $D(A)$ corresponding to $M|_U$ (Derived Categories of Schemes, Lemma
\ref{perfect-lemma-affine-compare-bounded})
has finite injective dimension\footnote{This condition is independent of the
choice of the affine open cover of the Noetherian scheme $X$.
Details omitted.}. To prove the lemma we may
replace $X$ by $U$, i.e., we may assume $X = \Spec(A)$
for some Noetherian ring $A$. Observe that
$R\SheafHom(K, L)$ is in $D^+_{\textit{Coh}}(\mathcal{O}_X)$ by
Derived Categories of Schemes, Lemma
\ref{perfect-lemma-coherent-internal-hom}.
Moreover, the formation of the left and right hand side
of the arrow commutes with the functor $D(A) \to D_\QCoh(\mathcal{O}_X)$ by
Lemma \ref{lemma-affine-duality} and
Derived Categories of Schemes, Lemma
\ref{perfect-lemma-quasi-coherence-internal-hom}
(to be sure this uses the assumptions on $K$, $L$, $M$ and what we just
proved about $R\SheafHom(K, L)$).
Then finally the arrow is an isomorphism by
More on Algebra, Lemmas
\ref{more-algebra-lemma-internal-hom-evaluate-isomorphism} part (2).

\medskip\noindent
Proof of (2). We argue as above. A small change is that here we get
$R\SheafHom(K, L)$ in $D_{\textit{Coh}}(\mathcal{O}_X)$ because
affine locally (which is allowable by Lemma \ref{lemma-affine-duality})
we may appeal to Dualizing Complexes, Lemma
\ref{dualizing-lemma-finite-ext-into-bounded-injective}.
Then we finally conclude by
More on Algebra, Lemma
\ref{more-algebra-lemma-internal-hom-evaluate-isomorphism-technical}.
\end{proof}

\end{document}
